\documentclass[12pt,a4paper]{article}
\usepackage[utf8]{inputenc}
\usepackage[T1]{fontenc}
\usepackage{geometry}
\geometry{margin=0.9in}
\usepackage{booktabs}
\usepackage{tabularx}
\usepackage{amsmath,amssymb}
\usepackage{hyperref}
\usepackage{xcolor}
\usepackage{enumitem}

\definecolor{primary}{HTML}{1B365D}
\definecolor{accent}{HTML}{2B6CB0}

\hypersetup{
    colorlinks=true,
    linkcolor=primary,
    urlcolor=accent,
    pdftitle={Automated Skin Disease Classification Using Deep Learning - Complete Official Report},
    pdfauthor={Anurag Saha, Ayan Karak}
}

\begin{document}

\begin{center}
    {\Large \textbf{\color{primary} JALPAIGURI GOVERNMENT ENGINEERING COLLEGE (JGEC)}}\\[0.3em]
    {\normalsize Department of Information Technology $\bullet$ 5th Semester Official Project Documentation}\\[1.2em]
    \rule{\linewidth}{1.5pt}\\[0.6em]
    {\LARGE \textbf{AUTOMATED SKIN DISEASE CLASSIFICATION USING DEEP LEARNING}}\\[0.4em]
    {\large \textit{Complete Official 5th Semester Project Report}}\\[0.4em]
    \rule{\linewidth}{1.5pt}\\[1em]
    \textbf{Submitted by:} Anurag Saha (Roll No: \texttt{24101106037}) \& Ayan Karak (Roll No: \texttt{24101106013})\\[0.3em]
    \textbf{Submitted to \& Supervised by:} Dr. Somnath Dhibar\\[0.3em]
    \textbf{Source Code Repository:} \url{https://github.com/sahaanurag207-oss/SKIN-DISEASE-CLASSIFICATION-FROM-IMAGES}
\end{center}

\vspace{1em}

\section{Executive Summary \& Project Overview}
Skin cancer and dermatological disorders represent a major global public health challenge, where early and accurate visual screening directly dictates clinical prognosis and survival rates. Traditional dermatoscopic examination relies heavily on specialized clinical expertise and suffers from substantial inter-observer variability due to high visual similarity across benign and malignant pigmented lesions. Developed as an official 5th Semester curriculum project in the Department of Information Technology at Jalpaiguri Government Engineering College (JGEC), this study presents an end-to-end, leakage-free deep learning framework for multi-class automated skin lesion classification.

The system is trained and rigorously evaluated on the benchmark \textbf{HAM10000} (``Human Against Machine with 10000 training images'') dataset curated by Tschandl et al., encompassing 10,015 dermatoscopic images distributed across seven distinct diagnostic categories. To overcome two pervasive pitfalls in medical computer vision---extreme class imbalance and patient-level data leakage---the pipeline integrates \texttt{StratifiedGroupKFold} partitioning grouped strictly by unique lesion identifiers (\texttt{lesion\_id}), a modernized ConvNeXt-Small backbone (\texttt{convnext\_small.fb\_in22k\_ft\_in1k}) pretrained on ImageNet-22k, domain-specific geometric and photometric Albumentations, Label Smoothing Cross-Entropy regularization, and multi-view Test-Time Augmentation (TTA).

\subsection*{Core Experimental Achievement}
On the strictly unseen, lesion-disjoint test partition ($n = 1,023$ images), the proposed ConvNeXt-Small + TTA pipeline achieves an \textbf{Overall Test Accuracy of 87.98\%}, a \textbf{Weighted Average F1-Score of 0.8761}, a \textbf{Macro Average F1-Score of 0.7534}, and a \textbf{Balanced Accuracy of 74.76\%}, notably reaching \textbf{97.44\% sensitivity (recall)} on Basal Cell Carcinoma (\texttt{bcc}) and \textbf{95.65\% F1-score} on Melanocytic Nevi (\texttt{nv}).

\subsection{Project Metadata \& Academic Identification}
\begin{table}[h!]
\centering
\renewcommand{\arraystretch}{1.25}
\begin{tabularx}{\textwidth}{>{\bfseries}l X}
\toprule
Attribute & Specification \\
\midrule
Project Title & Automated Skin Disease Classification Using Deep Learning \\
Institution & Jalpaiguri Government Engineering College (JGEC), West Bengal, India \\
Department \& Semester & Department of Information Technology --- 5th Semester Curriculum Project \\
Student Investigators & Anurag Saha (Roll No: 24101106037) \& Ayan Karak (Roll No: 24101106013) \\
Project Supervisor & Dr. Somnath Dhibar \\
Primary Dataset & HAM10000 (10,015 Dermatoscopic Images across 7 Diagnostic Classes) \\
Core Architecture & ConvNeXt-Small (\texttt{convnext\_small.fb\_in22k\_ft\_in1k}) via PyTorch \& \texttt{timm} \\
Validation Strategy & Leak-Free \texttt{StratifiedGroupKFold} Grouped by \texttt{lesion\_id} \\
Repository Source & \url{https://github.com/sahaanurag207-oss/SKIN-DISEASE-CLASSIFICATION-FROM-IMAGES} \\
\bottomrule
\end{tabularx}
\end{table}

\subsection{Target Dermatological Diagnostic Taxonomy (HAM10000)}
\begin{table}[h!]
\centering
\renewcommand{\arraystretch}{1.2}
\begin{tabularx}{\textwidth}{l l X r r}
\toprule
\textbf{Code} & \textbf{Diagnostic Category} & \textbf{Clinical Nature} & \textbf{Total} & \textbf{Test} \\
\midrule
\texttt{nv} & Melanocytic nevi & Benign melanocytic proliferation & 6,705 (66.95\%) & 684 \\
\texttt{mel} & Melanoma & Malignant melanocytic neoplasm (high lethality) & 1,113 (11.11\%) & 105 \\
\texttt{bkl} & Benign keratosis-like lesions & Benign (seborrheic keratoses, solar lentigo) & 1,099 (10.97\%) & 130 \\
\texttt{bcc} & Basal cell carcinoma & Malignant epithelial skin cancer & 514 (5.13\%) & 39 \\
\texttt{akiec} & Actinic keratoses \& IEC & Pre-malignant / in-situ squamous carcinoma & 327 (3.27\%) & 28 \\
\texttt{vasc} & Vascular lesions & Benign/reactive angiomas \& angiokeratomas & 142 (1.42\%) & 16 \\
\texttt{df} & Dermatofibroma & Benign fibrous cutaneous nodule & 115 (1.15\%) & 21 \\
\bottomrule
\end{tabularx}
\end{table}

\section{Literature Survey \& State-of-the-Art Review}
Automated dermatoscopic image analysis has evolved rapidly over the past decade, progressing from hand-crafted ABCD rule feature extractors to deep convolutional neural networks (CNNs), hierarchical Vision Transformers (ViTs), and modernized pure-convolutional architectures.

\subsection{Evolution of Deep Learning in Dermatoscopy}
\begin{itemize}[leftmargin=*]
    \item \textbf{Landmark CNN Feasibility (Esteva et al., Nature, 2017):} Demonstrated dermatologist-level skin cancer classification using an Inception-v3 architecture trained on 129,450 clinical and dermatoscopic images. While proving deep neural networks can match board-certified specialists, reliance on proprietary data highlighted the urgent need for standardized public benchmarks.
    \item \textbf{Standardization via HAM10000 (Tschandl et al., Scientific Data, 2018):} Released 10,015 multi-source dermatoscopic images across seven pigmented lesion categories with over 50\% histopathologic confirmation, establishing the \texttt{lesion\_id} field to track multiple views of the same physical lesion.
    \item \textbf{ISIC Grand Challenge Ensembles (Codella et al., 2018; Gessert et al., 2020):} Established that large-scale transfer learning, heavy geometric/color augmentation, label smoothing, and Test-Time Augmentation (TTA) consistently outperform single-crop inference on imbalanced dermatoscopy data.
    \item \textbf{Modernized Pure ConvNets --- ConvNeXt (Liu et al., CVPR, 2022):} Introduced $7\times 7$ depthwise convolutions, inverted bottlenecks, LayerNorm, and GELU activations, matching Vision Transformer accuracy while preserving convolutional inductive biases crucial for medical datasets of moderate scale ($\sim$10k images).
    \item \textbf{Lesion-Level Data Leakage Audits:} Reproducibility audits revealed that naive random or image-level stratified splits artificially inflate HAM10000 accuracy above 93--96\% because duplicate views of the same lesion cross between training and test partitions. Enforcing \texttt{lesion\_id} group partitioning eliminates this leakage.
\end{itemize}

\subsection{Comparative Synthesis of Prior Approaches vs. Proposed Framework}
\begin{table}[h!]
\centering
\renewcommand{\arraystretch}{1.25}
\begin{tabularx}{\textwidth}{l l l X}
\toprule
\textbf{Study / Paradigm} & \textbf{Backbone / Method} & \textbf{Split Protocol} & \textbf{Key Strengths \& Limitations} \\
\midrule
Traditional CAD (pre-2016) & Hand-crafted ABCD features + SVM & Random / Private sets & Interpretable, but fragile under lighting shifts, hair artifacts, and multi-class overlap. \\
Esteva et al. (2017) & Inception-v3 transfer learning & Clinician-curated split & Matched dermatologist sensitivity; required $>$129k private images. \\
Standard ResNet / DenseNet & ResNet-50 / DenseNet-121 + CE & Naive image-level split & Prone to severe lesion-level data leakage across duplicate HAM10000 views. \\
Vision Transformers (ViT/Swin) & Patch-based self-attention & 5-Fold Cross-Validation & Strong global context, but lacks translation equivariance on 10k images. \\
\textbf{Proposed JGEC Pipeline} & \textbf{ConvNeXt-Small (IN-22k) + LS + TTA} & \textbf{Strict \texttt{lesion\_id} disjoint} & \textbf{Zero lesion leakage; combines ViT macro design with CNN inductive bias.} \\
\bottomrule
\end{tabularx}
\end{table}

\section{Proposed Methodology \& System Architecture}
The proposed system implements a modular, end-to-end deep learning workflow designed specifically to address visual intra-class variability, severe class skewness (58:1 ratio between majority \texttt{nv} and minority \texttt{df} classes), and strict patient/lesion-level generalization.

\subsection{End-to-End Pipeline Architecture}
\begin{table}[h!]
\centering
\renewcommand{\arraystretch}{1.25}
\begin{tabularx}{\textwidth}{l l X}
\toprule
\textbf{Pipeline Stage} & \textbf{Component / Algorithm} & \textbf{Technical Configuration \& Purpose} \\
\midrule
1. Leakage Control & \texttt{StratifiedGroupKFold} & Groups all 10,015 images by unique \texttt{lesion\_id} while preserving class proportions; guarantees zero physical lesion overlap (Test $n = 1,023$). \\
2. Augmentation & Albumentations \& OpenCV & Stochastic \texttt{HorizontalFlip}, \texttt{VerticalFlip}, \texttt{RandomRotate90}, \texttt{ShiftScaleRotate}, and \texttt{ColorJitter}. \\
3. Backbone & ConvNeXt-Small (\texttt{timm}) & \texttt{convnext\_small.fb\_in22k\_ft\_in1k} leveraging $7\times 7$ depthwise convolutions for multi-scale pigment network extraction. \\
4. Classifier Head & Dropout + Linear Layer & Applies Dropout ($p = 0.3$) before the 7-class linear projection head to mitigate overfitting on minority classes. \\
5. Loss \& Optimizer & Label Smoothing CE + AdamW & Label smoothing ($\epsilon = 0.1$) with decoupled AdamW and \texttt{CosineAnnealingLR} decay schedule. \\
6. Acceleration & PyTorch AMP & FP16 mixed precision via \texttt{torch.amp.autocast} and \texttt{GradScaler} for fast GPU convergence. \\
7. Inference & Multi-View TTA & Averages softmax probabilities across original, horizontally flipped, and vertically flipped test views. \\
\bottomrule
\end{tabularx}
\end{table}

\subsection{Mathematical Formulation of Core Components}
\begin{enumerate}[leftmargin=*]
    \item \textbf{Leak-Free Lesion Partitioning:} Let $\mathcal{G}$ denote the set of unique lesion identifiers (\texttt{lesion\_id}) in HAM10000 and $g(x)$ map image $x$ to its physical lesion. \texttt{StratifiedGroupKFold} constructs disjoint lesion subsets $\mathcal{G}_{\text{train}}, \mathcal{G}_{\text{val}}, \mathcal{G}_{\text{test}}$ such that $\mathcal{G}_{\text{train}} \cap \mathcal{G}_{\text{test}} = \emptyset$ and $\mathcal{G}_{\text{val}} \cap \mathcal{G}_{\text{test}} = \emptyset$.
    \item \textbf{Label Smoothing Cross-Entropy Objective:} Given ground-truth class $y \in \{1, \dots, K\}$ ($K = 7$) and smoothing parameter $\epsilon = 0.1$, smoothed target probabilities are defined as:
    \begin{equation}
        q_k = \begin{cases} 1 - \epsilon, & \text{if } k = y \\ \dfrac{\epsilon}{K - 1}, & \text{if } k \neq y \end{cases}, \qquad \mathcal{L}_{\text{LS}} = -\sum_{k=1}^{K} q_k \log(p_k)
    \end{equation}
    \item \textbf{Test-Time Augmentation (TTA) Aggregation:} For each test sample $x_i$ and deterministic geometric transform set $\mathcal{T} = \{\text{Identity}, \text{HorizontalFlip}, \text{VerticalFlip}\}$, the final posterior probability distribution is computed as:
    \begin{equation}
        P_{\text{final}}(y \mid x_i) = \frac{1}{|\mathcal{T}|} \sum_{t \in \mathcal{T}} \text{Softmax}\big(f_\theta(t(x_i))\big)
    \end{equation}
\end{enumerate}

\section{Performance Chart \& Quantitative Analysis}
Model evaluation was conducted exclusively on the held-out, lesion-disjoint test partition comprising $n = 1,023$ dermatoscopic images across all seven diagnostic categories.

\subsection{Overall Aggregate Evaluation Metrics}
\begin{table}[h!]
\centering
\renewcommand{\arraystretch}{1.2}
\begin{tabularx}{\textwidth}{l c X}
\toprule
\textbf{Evaluation Metric} & \textbf{Score} & \textbf{Analytical Interpretation} \\
\midrule
\textbf{Overall Test Accuracy} & \textbf{87.98\%} & High global classification correctness across 1,023 unseen lesion images. \\
\textbf{Weighted Average F1-Score} & \textbf{0.8761} & Strong precision-recall harmonic balance proportional to clinical prevalence. \\
\textbf{Macro Average F1-Score} & \textbf{0.7534} & Robust unweighted performance across both majority and rare ($<2\%$) classes. \\
\textbf{Balanced Accuracy} & \textbf{74.76\%} & Mean per-class recall across all 7 categories despite a 58:1 dataset imbalance. \\
\bottomrule
\end{tabularx}
\end{table}

\subsection{Per-Class Precision, Recall, and F1-Score Breakdown}
\begin{table}[h!]
\centering
\renewcommand{\arraystretch}{1.2}
\begin{tabularx}{\textwidth}{l X c c c r}
\toprule
\textbf{Class} & \textbf{Diagnostic Name} & \textbf{Precision} & \textbf{Recall} & \textbf{F1-Score} & \textbf{Support} \\
\midrule
\texttt{nv} & Melanocytic nevi & 0.9483 & 0.9649 & 0.9565 & 684 \\
\texttt{mel} & Melanoma & 0.7160 & 0.5524 & 0.6237 & 105 \\
\texttt{bkl} & Benign keratosis-like lesions & 0.7518 & 0.7923 & 0.7715 & 130 \\
\texttt{bcc} & Basal cell carcinoma & 0.6786 & 0.9744 & 0.8000 & 39 \\
\texttt{akiec} & Actinic keratoses \& intraepithelial carcinoma & 0.6667 & 0.5000 & 0.5714 & 28 \\
\texttt{vasc} & Vascular lesions & 0.7333 & 0.6875 & 0.7097 & 16 \\
\texttt{df} & Dermatofibroma & 0.9412 & 0.7619 & 0.8421 & 21 \\
\bottomrule
\end{tabularx}
\end{table}

\subsection{Diagnostic Error \& Class Imbalance Insights}
\begin{itemize}[leftmargin=*]
    \item \textbf{High-Sensitivity Malignancy Screening (\texttt{bcc}):} Basal cell carcinoma achieved an exceptional Recall of 97.44\% ($\text{F1} = 0.8000$) despite having only 39 test samples, demonstrating that ConvNeXt's $7\times 7$ depthwise receptive fields effectively capture telangiectatic vessels and pearly borders.
    \item \textbf{Distinct Morphological Recognition in Rare Classes (\texttt{df} \& \texttt{vasc}):} Dermatofibroma ($n = 21$) achieved 94.12\% Precision and 0.8421 F1-Score, while Vascular lesions ($n = 16$) reached 0.7097 F1-Score, confirming strong feature transfer from ImageNet-22k weights under extreme sample scarcity.
    \item \textbf{Melanoma vs. Nevi Visual Overlap (\texttt{mel}):} Melanoma recorded 71.60\% Precision and 55.24\% Recall ($\text{F1} = 0.6237$). Early-stage melanoma shares significant pigment network symmetry with atypical nevi (\texttt{nv}), causing borderline cases to default toward the dominant \texttt{nv} prior (66.95\% of dataset).
    \item \textbf{Pre-Malignant Boundary Ambiguity (\texttt{akiec}):} Actinic keratoses ($n = 28$) yielded an F1-Score of 0.5714 (Precision 66.67\%, Recall 50.00\%), attributable to subtle scaly surface textures overlapping with benign keratosis-like lesions (\texttt{bkl}, $\text{F1} = 0.7715$).
\end{itemize}

\section{Conclusion, Acknowledgment \& Future Scope}
\subsection{Conclusion}
This project successfully designs, implements, and validates an automated deep learning pipeline for seven-class dermatoscopic skin disease classification using the HAM10000 benchmark. By replacing naive image-level splitting with lesion-grouped \texttt{StratifiedGroupKFold} validation, the study establishes an honest, leakage-free evaluation protocol suitable for clinical decision-support research. Leveraging a pretrained ConvNeXt-Small backbone combined with domain-tailored Albumentations, Label Smoothing Cross-Entropy, AdamW optimization, PyTorch AMP, and multi-view TTA, the model attains 87.98\% Overall Test Accuracy, 0.8761 Weighted F1-Score, and 0.7534 Macro F1-Score.

\subsection{Acknowledgment}
We express our sincere gratitude and appreciation to \textbf{Dr. Somnath Dhibar}, Department of Information Technology, Jalpaiguri Government Engineering College (JGEC), for his invaluable mentorship, technical guidance, and continuous support throughout the design, validation, and development of this project.

\subsection{Limitations \& Future Scope of Research}
\begin{enumerate}[leftmargin=*]
    \item \textbf{Cost-Sensitive \& Focal Loss Optimization:} Integrating Class-Balanced Focal Loss, LDAM loss, and post-hoc ROC threshold tuning to elevate clinical recall on high-risk malignant classes (\texttt{mel} and \texttt{akiec}).
    \item \textbf{Multimodal Patient Metadata Fusion:} Fusing clinical covariates (patient age, sex, and anatomical lesion site) with ConvNeXt visual embeddings via cross-attention layers to disambiguate visually similar lesions.
    \item \textbf{Explainable AI (XAI) for Clinical Trust:} Integrating Grad-CAM++ and Eigen-CAM visual saliency overlays to highlight dermatoscopic structures driving predictions.
    \item \textbf{Multi-Model Ensembling \& Edge Deployment:} Ensembling ConvNeXt-Small with Swin Transformer and EfficientNetV2 backbones via ONNX/TensorRT quantization for real-time teledermatology screening.
\end{enumerate}

\end{document}
